\begin{table}[H]
\centering
\caption{How far the unmatched would have to depart from missing-at-random}
\label{tab:tipping}
\begin{threeparttable}
\begin{tabular}{cc}
\toprule
$\delta$ (log-odds shift) & Mid-pandemic adjusted gap \\
\midrule
0.0000 & 0.0186 \\
-0.2500 & 0.0175 \\
-0.5000 & 0.0166 \\
-1.0000 & 0.0153 \\
-2.0000 & 0.0140 \\
-4.0000 & 0.0132 \\
-8.0000 & 0.0131 \\
\midrule
Tipping point & none in $\delta \geq -8$ \\
Gap in the limit & 0.0131 \\
\bottomrule
\end{tabular}
\begin{tablenotes}[flushleft]\footnotesize
\item Notes: Each row re-estimates equation~\eqref{eq:es} on all origins for which a further interview is scheduled, with the outcome of unmatched workers imputed from a model fitted on matched workers of the same education group and then shifted by $\delta$ in log-odds for graduates only. $\delta = 0$ is missing-at-random within education group; negative values make unmatched graduates leave employment less often than their observed characteristics imply, the direction that works against the reversal. The tipping point is the value at which the mid-pandemic adjusted gap reaches zero, located by bisection. Because only the difference between the two groups' shifts matters, the non-graduate shift is fixed at zero without loss.
\end{tablenotes}
\end{threeparttable}
\end{table}
